% id: 201-12
% book: 베이직쎈 대수 · 12 수학적 귀납법
% section: 유형 124 수열의 귀납적 정의의 활용
% figure: none
% answer: 
% answer_source: 
% variant_level: 0 (원본 전사)
% 이 파일은 build.mjs 가 items.json 에서 만든다. 고칠 때는 items.json 을 고친다.
\smash{\raisebox{1.5mm}{\dmkichul{베이직쎈 대수 201쪽 12번}}}
어느 여행 동아리 모임에 참석한 $n$명의 회원이 모두 서로 한 번씩 악수를 하였다. 다음은 $n$명의 회원이 악수한 총횟수를 $a_n$이라 할 때, $a_1$을 구하고, $a_n$과 $a_{n+1}$ 사이의 관계식을 구하는 과정이다.\exprbox{\begin{minipage}{0.96\linewidth}\raggedright 한 명의 회원이 참석하면 악수를 할 수 없으므로\\ \hspace*{1.5em}$a_1=\fbox{㈎}$\\ $n$명의 회원이 악수한 총횟수는 $a_n$이고, $n$명의 회원이 모두 악수를 한 후 $(n+1)$번째 회원이 참석하면 $(n+1)$번째 회원은 먼저 와 있던 $n$명의 회원과 한 번씩 악수를 하게 되므로\\ \hspace*{1.5em}$a_{n+1}=a_n+\fbox{㈏}$ $(n=1{,}\mkern3mu\mathopen{}2{,}\mkern3mu\mathopen{}3{,}\mkern3mu\mathopen{}\ldots)$\end{minipage}}위의 과정에서 ㈎, ㈏에 알맞은 것을 구하시오.
